\section{Conclusion}\label{sec:concl}
We have shown that the APEIRON Layer~2--3 pipeline — originally validated on image data — transfers without modification to synthetic benchmarks for electrocardiograms, financial time series, and seismic records.
In each domain, it autonomously induces functional modules that correspond to clinically, economically, or geophysically meaningful structures.
This confirms that the APEIRON substrate is a genuinely domain‑agnostic engine for unsupervised knowledge genesis.

From the perspective of the seventeen‑layer architecture, these experiments demonstrate the first, crucial phase of \emph{Symmetry Breaking}: the spontaneous emergence of distinguishable categories from an initially uniform pool of observables.
The fact that this emergence occurs autonomously, without human labels or feature engineering, validates the core premise of the APEIRON Framework: that intelligence can be bootstrapped from first principles, directly from the relational structure of elementary data.

Future work will extend the pipeline to higher layers (Layer~4 and beyond) to capture temporal dynamics and multi‑scale interactions, and will apply the full seventeen‑layer architecture to real‑world autonomous monitoring problems in healthcare, finance, and geophysics.
A particularly important next step is the \emph{autonomous calibration} of the window size and co‑activation threshold, so that the system can self‑tune to the characteristic time‑scales of each domain — a requirement for genuine non‑anthropocentric intelligence.